% !TEX program = lualatex
\documentclass[11pt,a4paper]{article}
\usepackage{luatexja}
\usepackage{amsmath,amssymb,amsthm,amsfonts}
\usepackage{geometry}
\usepackage{enumitem}
\usepackage{hyperref}
\geometry{margin=1in}

%-------------------------------------------------
%  定理環境
%-------------------------------------------------
\theoremstyle{definition}
\newtheorem{definition}{定義}[section]
\newtheorem{axiom}{公理}[section]
\newtheorem{lemma}{補題}[section]
\newtheorem{theorem}{定理}[section]
\newtheorem{corollary}{系}[section]
\newtheorem{proposition}{命題}[section]
\newtheorem{remark}{注記}[section]
\newtheorem{example}{例}[section]

\renewcommand{\proofname}{\textbf{証明.}}
\renewcommand{\abstractname}{Abstract}

%-------------------------------------------------
%  マクロ
%-------------------------------------------------
\newcommand{\Mr}{\mathcal{M}_{\mathrm{r}}}
\newcommand{\etar}{\eta^{\mathrm{r}}}
\newcommand{\mur}{\mu^{\mathrm{r}}}
\newcommand{\id}{\mathrm{id}}
\newcommand{\Hom}{\mathrm{Hom}}
\newcommand{\End}{\mathbf{End}}
\DeclareMathOperator*{\bigsqcap}{\sqcap}
\providecommand{\Ref}{\mathbf{Ref}}
\newcommand{\Dom}{\mathbf{Dom}}

%-------------------------------------------------
\begin{document}

\title{洗練モナド：\\
洗練順序領域理論における前進閉包作用素の定式化}
\author{千葉将臣}
\date{2026-06-29（2026-07-02改訂）}
\maketitle

\begin{abstract}
本稿は洗練モナド $(\Mr, \etar, \mur)$ を定式化する。
洗練モナドは、対象を洗練文脈へ持ち上げ、洗練の洗練を一段階へ統合する
前進方向の閉包作用素である。
従来稿では洗練コモナドおよびフロベニウス条件を洗練構造そのものに含めていたが、
本稿ではそれらを基本定義から除外する。
フロベニウス性は、搾取モナド、二重否定モナド、洗練モナド、カット除去モナドに
統制モナドを加えた5モナド構成において、識軸スパイラルから誘導される性質として扱う。
したがって、洗練モナド自体は通常のモナドであり、
その役割は洗練連鎖の進行、文脈化、合成、極限収束を記述することに限定される。
\end{abstract}

\tableofcontents
\newpage

%=============================================================
\section{序論：洗練の前進構造}
%=============================================================

\subsection{通常のモナドとしての洗練}

Moggi \cite{Moggi1991} によるモナドを用いた計算の意味論以来、
副作用・状態・継続といった計算の文脈はモナドの
単位 $\eta$ と乗法 $\mu$ によって記述されてきた。
本稿では、洗練連鎖も同様に通常のモナドとして定式化する。

洗練モナド $\Mr$ は、対象 $A$ を「洗練過程の中にある $A$」へ持ち上げる。
単位 $\etar_A:A\to\Mr A$ は対象の文脈化を表し、
乗法 $\mur_A:\Mr\Mr A\to\Mr A$ は洗練の洗練を一つの洗練過程へ統合する。

この構成において、洗練モナドは完了確定の逆方向を基本構造として持たない。
完了確定や後退方向は、別稿で扱う5モナド構成、すなわち統制モナドを識軸とする
スパイラル構成から誘導される。

\subsection{洗練順序領域理論の動機}

Abramsky \cite{Abramsky1991} のDomain Theory in Logical Form は
直観主義論理と領域理論を接続したが、
情報順序 $\sqsubseteq$ における $\bot$（停止しない計算）を最小元として扱う。

\begin{equation}
\bot \sqsubseteq a \sqsubseteq \top \quad \text{（情報順序）}
\end{equation}

これに対し、洗練順序では $\bot$ を未確定な全可能性として扱う。
完了した記述 $a$ は、全可能性 $\bot$ から可能性を削減して得られる。

\begin{equation}
\bot \sqsupseteq a \sqsupseteq \top \quad \text{（洗練順序）}
\end{equation}

本稿は洗練順序を採用し、$\bot$ をRefinement Orderにおける最大元として扱う。
この反転により、計算の進行は順序の降下として記述される。
Hoare と He \cite{Hoare1998} のUnifying Theories of Programming（UTP）における
洗練順序と情報順序の区別を本稿の出発点とする。

\subsection{本稿の目的}

本稿は以下を達成する。

\begin{enumerate}
\item 洗練モナド $(\Mr,\etar,\mur)$ を圏論的に定式化する。
\item 洗練順序領域における洗練連鎖とモナド構造の整合性を示す。
\item 洗練連鎖の極限としての不動点収束を記述する。
\item 他のモナドとの合成において生じる非可換性を洗練残余として定義する。
\item フロベニウス性は洗練モナド単体の前提ではなく、5モナド構成で担保されることを明確化する。
\end{enumerate}

%=============================================================
\section{予備的概念}
%=============================================================

\subsection{洗練順序と領域}

\begin{definition}[洗練順序領域]
半順序集合 $(D, \sqsupseteq)$ が洗練順序領域であるとは、
以下を満たすことをいう。

\begin{enumerate}
\item $D$ は最大元 $\bot \in D$（全可能性）を持つ。
\item $D$ は洗練連鎖の下限を持つ。
\item 任意の有向集合 $S \subseteq D$ に対し、
  $\bigsqcap S$ が存在する。
\end{enumerate}

$a \sqsupseteq b$ は「$a$ は $b$ より未確定であり、より多くの可能性を含む」ことを意味する。
\end{definition}

\begin{remark}
情報順序では $\bot$ が最小元（情報なし）であるのに対し、
洗練順序では $\bot$ が最大元（全可能性）である。
この反転により、計算の進行は順序の降下として記述される。
\end{remark}

\begin{definition}[洗練連鎖]
洗練順序領域 $(D, \sqsupseteq)$ における洗練連鎖とは、
降下列
\begin{equation}
\bot = a_0 \sqsupseteq a_1 \sqsupseteq a_2 \sqsupseteq \cdots
\end{equation}
であって、その下限 $\bigsqcap_n a_n$ が存在するものをいう。
\end{definition}

\subsection{自己関手圏}

\begin{definition}[自己関手圏]
圏 $\mathcal{C}$ の自己関手圏を
\begin{equation}
\End(\mathcal{C}) = [\mathcal{C},\mathcal{C}]
\end{equation}
と書く。
対象は自己関手 $F:\mathcal{C}\to\mathcal{C}$、射は自然変換である。
関手合成をモノイダル積、恒等関手を単位対象とみなす。
\end{definition}

%=============================================================
\section{洗練モナドの定式化}
%=============================================================

\subsection{洗練文脈}

\begin{definition}[洗練文脈 $\Mr$]
洗練文脈 $\Mr$ を、対象が洗練順序において
より確定した記述へ移行する過程を包摂する文脈として定義する。
$\Mr A$ は「$A$ が洗練過程の中にある状態」を表す。
\end{definition}

\begin{definition}[洗練モナド]
圏 $\mathcal{C}$ 上の自己関手 $\Mr:\mathcal{C}\to\mathcal{C}$ と
自然変換 $\etar:\id_{\mathcal{C}}\Rightarrow\Mr$、
$\mur:\Mr\Mr\Rightarrow\Mr$ からなる三つ組
\begin{equation}
(\Mr,\etar,\mur)
\end{equation}
を洗練モナドとして定義する。

\begin{align}
\etar_A &: A \to \Mr A
  \quad \text{（対象を洗練文脈に持ち上げる）} \\
\mur_A &: \Mr\Mr A \to \Mr A
  \quad \text{（洗練の洗練を一段階に統合する）}
\end{align}

モナド則は以下である。

\begin{align}
\mur \circ \Mr\mur &= \mur \circ \mur\Mr
  \quad \text{（結合則）} \\
\mur \circ \Mr\etar &= \mur \circ \etar\Mr = \id_{\Mr}
  \quad \text{（単位元則）}
\end{align}
\end{definition}

\begin{proposition}[洗練順序との整合性]
洗練モナドにおいて、$\etar_A:A\to\Mr A$ は
対象を洗練過程へ導入する操作として定式化される。
また、$\Mr$ は洗練順序に関して単調である。

\begin{equation}
a \sqsupseteq a' \Rightarrow \Mr a \sqsupseteq \Mr a'
\end{equation}

洗練連鎖 $\bot \sqsupseteq a_1 \sqsupseteq a_2 \sqsupseteq \cdots$ に対し、
$\Mr$ の適用は洗練過程を一段持ち上げた文脈を生成する。
\end{proposition}

\begin{proof}
$\Mr$ は自己関手であるため、対象と射を保つ。
洗練順序を圏 $\mathcal{C}$ の射として読むと、
$a\sqsupseteq a'$ は $a$ から $a'$ への洗練射を与える。
関手性により、この射は $\Mr a$ から $\Mr a'$ への射に写される。
したがって $\Mr$ は洗練順序に関して単調である。
\end{proof}

\subsection{洗練合成}

\begin{definition}[洗練合成]
$\mur_A:\Mr\Mr A\to\Mr A$ は、洗練された対象をさらに洗練する二段階過程を、
一つの洗練過程へ圧縮する操作である。
\begin{equation}
\Mr(\Mr A) \xrightarrow{\mur_A} \Mr A
\end{equation}
\end{definition}

\begin{remark}
洗練合成は、仕様から設計、設計から実装へ進むような多段階の洗練を、
単一の洗練連鎖として扱うための構造である。
これはUTPにおける仕様・設計・実装の洗練列と対応する。
\end{remark}

%=============================================================
\section{不動点収束}
%=============================================================

\begin{definition}[洗練不動点]
対象 $A$ に対し、洗練モナドの反復列
\begin{equation}
A \xrightarrow{\etar_A} \Mr A \xrightarrow{\Mr\etar_A} \Mr^2 A
\xrightarrow{\Mr^2\etar_A} \cdots
\end{equation}
が洗練順序上で極限を持つとき、その極限を洗練不動点 $A^*$ と呼ぶ。
\begin{equation}
A^* = \bigsqcap_{n\geq 0}\Mr^n A
\end{equation}
\end{definition}

\begin{theorem}[洗練連鎖の収束]
洗練順序領域 $(D,\sqsupseteq)$ が洗練連鎖の下限を持ち、
$\Mr$ が連続な自己関手であるならば、任意の対象 $A$ に対して
洗練反復列 $\{\Mr^n A\}_{n\geq 0}$ は下限 $A^*$ を持つ。
\end{theorem}

\begin{proof}
仮定より、$\{\Mr^n A\}_{n\geq 0}$ は洗練順序上の連鎖である。
洗練順序領域は連鎖の下限を持つため、
$\bigsqcap_{n\geq 0}\Mr^n A$ が存在する。
これを $A^*$ と定義する。
$\Mr$ の連続性により、$\Mr$ はこの下限を保存するため、
$A^*$ は洗練反復に対する不動点として振る舞う。
\end{proof}

\begin{corollary}[$\bot$ との同型性]
洗練順序における $\bot$（全可能性・最大元）は、
洗練連鎖の出発点であると同時に、不動点収束を記述するための基準点である。

\begin{equation}
\bot \leadsto A^* = \bigsqcap_{n\geq 0}\Mr^n A
\end{equation}

情報順序における欠如としての $\bot$ ではなく、
洗練順序における全可能性としての $\bot$ が、
洗練過程の開始条件を与える。
\end{corollary}

%=============================================================
\section{洗練残余}
\label{sec:residue}
%=============================================================

\begin{definition}[洗練残余]
洗練モナド $\Mr$ と別のモナド $\mathcal{M}$ の合成が可換でない領域を、
洗練残余 $\mathcal{R}(\Mr,\mathcal{M})$ として定義する。

\begin{equation}
\Mr(\mathcal{M}A) \not\cong \mathcal{M}(\Mr A)
\quad\Rightarrow\quad
\mathcal{R}(\Mr,\mathcal{M}) \neq \emptyset
\end{equation}
\end{definition}

\begin{theorem}[残余の確定]
$\Mr\circ\mathcal{M}$ と $\mathcal{M}\circ\Mr$ が同型でないとき、
二つの洗練経路の差分は洗練残余として確定する。
\end{theorem}

\begin{proof}
自然変換 $\theta:\Mr\mathcal{M}\Rightarrow\mathcal{M}\Mr$ が
同型でないとき、二つの合成関手は同一の洗練過程を表さない。
この非同型性は、$\Mr$ による文脈化と $\mathcal{M}$ による文脈化の順序に依存する差分を生む。
その差分を $\mathcal{R}(\Mr,\mathcal{M})$ と呼ぶ。
直観主義論理において $\neg\neg P\to P$ が一般に成立しないことと同様に、
合成順序の非可換性は単なる失敗ではなく、残余として正の対象化を受ける。
\end{proof}

\begin{remark}[継続残余との対応]
継続残余 $\rho$（千葉 \cite{Chiba2026continuation}）の語彙では、
$\Mr\mathcal{M}\cong\mathcal{M}\Mr$ の場合は $\rho=0$ に対応し、
非可換な場合は $\rho>0$ に対応する。
洗練残余は、この $\rho>0$ を圏論的な合成順序の差分として記述する。
\end{remark}

\subsection{カット除去モナドとの関係}

\begin{proposition}[カット除去モナドとの共存]
カット除去モナド $\mathcal{C}$（千葉 \cite{Chiba2026cutmonad}）は、
洗練モナド $\Mr$ の特殊ケースではなく、洗練連鎖と共存する独立した通常モナドである。

\begin{equation}
\Mr\circ\mathcal{C}\cong\mathcal{C}\circ\Mr
\end{equation}
が成立する場合、洗練の前進とカット除去の前進は可換に共存する。
成立しない場合、その差分は $\mathcal{R}(\Mr,\mathcal{C})$ として確定する。
\end{proposition}

\begin{proof}
カット除去モナド $\mathcal{C}$ は証明図を正規化する前進方向のモナドである。
一方、洗練モナド $\Mr$ は対象を洗練文脈へ持ち上げる。
両者は対象とする層が異なるため、一方が他方を包摂するのではなく、
合成可能な独立構造として扱われる。
合成が可換ならば両者は整合的に共存し、非可換ならばその差分が残余となる。
\end{proof}

%=============================================================
\section{5モナド構成との関係}
%=============================================================

\begin{proposition}[フロベニウス性の外部化]
洗練モナド $\Mr$ は、それ自体の定義に洗練コモナドやフロベニウス条件を含まない。
フロベニウス性は、搾取モナド、二重否定モナド、洗練モナド、カット除去モナド、
統制モナドからなる5モナド構成において担保される。
\end{proposition}

\begin{proof}
洗練モナドの基本データは $(\Mr,\etar,\mur)$ であり、
これは通常のモナドのデータと一致する。
したがって、洗練モナド単体からは余単位や余乗法は要求されない。
後退方向やフロベニウス性は、統制モナドを識軸とするスパイラル構成によって、
洗練モナドとカット除去モナドへ誘導される。
ゆえにフロベニウス性は洗練モナドの定義ではなく、
5モナド全体の構成的帰結である。
\end{proof}

\begin{remark}
この整理により、洗練モナドは前進方向の洗練を担う単純な構造として保持される。
洗練の完了確定、方向反転、双方向性は、統制モナドを含む上位構成で扱う。
したがって、洗練モナドと5モナド構成の役割分担が明確になる。
\end{remark}

%=============================================================
\section{既存理論との対比}
%=============================================================

\subsection{通常のモナドとの対比}

\begin{center}
\begin{tabular}{lll}
\hline
& \textbf{通常のモナド} & \textbf{洗練モナド} \\
\hline
方向性 & 前進のみ（$\eta,\mu$） & 前進のみ（$\etar,\mur$） \\
文脈 & 副作用・状態・継続など & 洗練過程 \\
完了条件 & 外部から固定される場合が多い & 5モナド構成で扱う \\
$\bot$ の扱い & 最小元（情報順序） & 全可能性（洗練順序） \\
残余 & 通常は明示されない & 合成非可換性として記述 \\
\hline
\end{tabular}
\end{center}

\subsection{二重否定モナドとの対比}

二重否定モナド $\neg\neg$ は、直観主義論理を古典論理へ反射する閉包作用素である。
一方、洗練モナドは古典化を目的としない。
洗練モナドは $\bot$ を排除せず、全可能性として保持しながら、
洗練連鎖の進行を記述する。

\begin{center}
\begin{tabular}{lll}
\hline
& \textbf{二重否定モナド} & \textbf{洗練モナド} \\
\hline
作用 & $A\mapsto\neg\neg A$ & $A\mapsto\Mr A$ \\
目的 & 古典的安定化 & 洗練過程の文脈化 \\
$\bot$ & 閉包外へ押し出す傾向 & 全可能性として保持 \\
残余 & $\neg\neg A\to A$ の非成立 & 合成非可換性 \\
\hline
\end{tabular}
\end{center}

%=============================================================
\section{適用域}
%=============================================================

洗練モナドはプログラミングに限定されない。
洗練連鎖として記述可能な任意の過程に適用できる。

\begin{center}
\begin{tabular}{ll}
\hline
\textbf{適用域} & \textbf{$\Mr A$ の内容} \\
\hline
プログラミング & コードの洗練過程 \\
学習 & 理解の深化過程 \\
対話 & 議論の精緻化過程 \\
芸術制作 & 表現の洗練過程 \\
科学研究 & 仮説の洗練過程 \\
\hline
\end{tabular}
\end{center}

%=============================================================
\section{結論}
%=============================================================

本稿は洗練モナド $(\Mr,\etar,\mur)$ を以下の観点から定式化した。

\begin{enumerate}
\item \textbf{理論的基盤}：
  洗練順序領域理論において $\bot$ を全可能性として扱う
  情報順序の反転を採用し、UTPの洗練順序・情報順序の区別を出発点とした。

\item \textbf{圏論的構成}：
  洗練モナド $(\Mr,\etar,\mur)$ を、対象を洗練文脈へ持ち上げ、
  洗練の洗練を一段階へ統合する通常のモナドとして定義した。

\item \textbf{不動点収束}：
  洗練連鎖の極限として不動点 $A^*$ への収束を記述した。

\item \textbf{残余の確定}：
  他のモナドとの合成非可換性を、洗練残余 $\mathcal{R}(\Mr,\mathcal{M})$ として定義した。

\item \textbf{5モナド構成との役割分担}：
  フロベニウス性を洗練モナド単体の定義から外し、
  統制モナドを含む5モナド構成によって担保される性質として位置づけた。
\end{enumerate}

洗練モナドは完成した双方向構造ではなく、洗練の前進方向を担う基本モナドである。
方向反転、完了確定、フロベニウス性は、統制モナドを識軸とする上位構成で扱われる。
この分離により、洗練モナドの定義は簡潔になり、5モナド構成における役割も明確になる。

\begin{thebibliography}{99}

\bibitem{Abramsky1991}
Samson Abramsky.
\newblock Domain theory in logical form.
\newblock \emph{Annals of Pure and Applied Logic}, 51(1--2):1--77, 1991.

\bibitem{Hoare1998}
C.~A.~R. Hoare and He Jifeng.
\newblock \emph{Unifying Theories of Programming}.
\newblock Prentice Hall, 1998.

\bibitem{MacLane1971}
Saunders Mac Lane.
\newblock \emph{Categories for the Working Mathematician}.
\newblock Springer, 1971.

\bibitem{Moggi1991}
Eugenio Moggi.
\newblock Notions of computation and monads.
\newblock \emph{Information and Computation}, 93(1):55--92, 1991.

\bibitem{Plotkin1977}
Gordon Plotkin.
\newblock LCF considered as a programming language.
\newblock \emph{Theoretical Computer Science}, 5(3):223--255, 1977.

\bibitem{Scott1970}
Dana Scott.
\newblock Outline of a mathematical theory of computation.
\newblock In \emph{Proceedings of the 4th Annual Princeton Conference
  on Information Sciences and Systems}, pages 169--176, 1970.

\bibitem{Chiba2026cutmonad}
千葉将臣.
\newblock カット消去モナド：証明構造の安定化を担うモナド的定式化.
\newblock 未刊行原稿, 2026.

\bibitem{Chiba2026continuation}
千葉将臣.
\newblock 継続残余としての証明障害：実行可能性と静的証明のあいだの初等モデル.
\newblock 未刊行原稿, 2026.

\end{thebibliography}

\end{document}
